%
% $Id: slides.tex 4228 2006-06-21 21:55:12Z jjamor $
%
%
% Compilar a .pdf con LaTeX (pdflatex)
% Es necesario instalar Beamer (paquete latex-beamer en Debian)
%

%
% Gráficos:
% Los gráficos pueden suministrarse en PNG, JPG, TIF, PDF, MPS
% Los EPS deben convertirse a PDF (usar epstopdf)
%

\documentclass{beamer}
\usetheme{Warsaw}
\usebackgroundtemplate{\includegraphics[width=\paperwidth]{format/libresoft-bg.png}}
\usepackage[spanish]{babel}
\usepackage[utf8]{inputenc}
\usepackage{graphics}
\usepackage{amssymb} % Simbolos matematicos
\usepackage{url}

%\definecolor{libresoftgreen}{RGB}{162,190,43}
%\definecolor{libresoftblue}{RGB}{0,98,143}

%\setbeamercolor{titlelike}{bg=libresoftgreen}

%% Metadatos del PDF.
\hypersetup{
  pdftitle={Introduction to software quality evaluation},
  pdfauthor={F. Ortega, D. Izquierdo, P. Coca},
  pdfcreator={GSyC/Libresoft},
  pdfproducer=PDFLaTeX,
  pdfsubject={nn},
}
%%


\AtBeginSection[]
{
  \begin{frame}<presentation>
    \frametitle{Index}
    \tableofcontents[current]
  \end{frame}
}


\begin{document}

\title{Introduction to software quality evaluation}
\subtitle{Master on Free Software}
\institute{\\jfelipe@libresoft.es\\
GSyC/Libresoft}
\author{Felipe Ortega, Daniel Izquierdo, Pedro Coca}
\date{\today}

\frame{
\maketitle
\begin{center}
\includegraphics[width=6cm]{format/gsyc-urjc}
\end{center}
}


% Si el titulo o el autor se quieren acortar para los pies de página
% se pueden redefinir aquí:
%\title{Titulo corto}
%\author{Autores abreviado}


%% LICENCIA DE REDISTRIBUCION DE LAS TRANSPAS
\frame{
~
\vspace{4cm}

\begin{flushright}
{\tiny
(cc) 2010 Felipe Ortega, Daniel Izquierdo, Pedro Coca. \\
Some rights reserved. This document is distributed under the Creative \\
            Commons Attribution-ShareAlike 3.0 licence, available in \\
            http://creativecommons.org/licenses/by-sa/3.0/

%  Este documento (o uno muy similar) está disponible en \\
%  \url{http://gsyc.escet.urjc.es/~jjamor/}
}
\end{flushright}
}
%%

%%%%%%
%Transpas separadas por \begin{frame}
%%%%%%%%%%%%%%%%%%%%%%%%\end{frame}

\section{Introduction}

\begin{frame}
\frametitle{What is software quality?}
\begin{itemize}
\item Complex concept, different from quality in manufacturing.
\begin{itemize}
 \item Product developed to meet its specification.
\end{itemize}

\item But it is difficult to write good software specification.
\item Maintainability, reusability, modularity, security, etc. 
are usually not included (but required).
\item Quality is linked to meeting end-user expectations (subjective).
\end{itemize}
\end{frame}

%%%%%%%%%%%%%%%%%%%%%%%%%%%%%%%%%%%%%%%%%%%%%%%%%%%%%%%%%%%%%%

\begin{frame}
 \frametitle{Factors that may improve quality}
 \begin{itemize}
 \item OOP, CASE support.
 \item Standards (development and organizational). But bureaucracy is not enough.
 \item Quality culture: experience, best practices, management.
 \item But they also find a common stratification pattern:
  \textbf{the onion model}
 \end{itemize}
\end{frame}

%%%%%%%%%%%%%%%%%%%%%%%%%%%%%%%%%%%%%%%%%%%%%%%%%%%%%%%%%%%%%%

\begin{frame}
 \frametitle{Main areas of software quality in large projects}
 \begin{itemize}
  \item \textbf{Quality assurance}: Defining framework for organizational quality.
  \item \textbf{Quality planning}: Selecting elements from the framework, adapted
  to a particular project.
  \item \textbf{Quality control}: Checking that development teams adhere to the plan.
  \item For many \textit{community} FLOSS projects, none of the above are defined.
  \item Problems for companies with this methods in place to integrate community FLOSS.
 \end{itemize}
\end{frame}

%%%%%%%%%%%%%%%%%%%%%%%%%%%%%%%%%%%%%%%%%%%%%%%%%%%%%%%%%%%%%%

\begin{frame}
 \frametitle{Process and product quality}
 \begin{itemize}
 \item In traditional manufacturing, process quality is tightly linked to final product quality.
 \item But software involves design, creativity. Subject to past experience and team skills.
 \item Difficult to measure some aspects: maintainability, reusability, etc.
 \item Obscure links between changes in process and software product quality.
 \item Blind application of process quaility is not good for all situations.
 \end{itemize}
\end{frame}

%%%%%%%%%%%%%%%%%%%%%%%%%%%%%%%%%%%%%%%%%%%%%%%%%%%%%%%%%%%%%%

\section{Quality Assurance (QA): standards}

%%%%%%%%%%%%%%%%%%%%%%%%%%%%%%%%%%%%%%%%%%%%%%%%%%%%%%%%%%%%%%

\begin{frame}
\frametitle{Quality assurance overview}
\begin{center}
\begin{figure}
 \includegraphics[height=7.5cm]{figs/quality-assurance.pdf}
\end{figure}
\end{center}
\end{frame}

%%%%%%%%%%%%%%%%%%%%%%%%%%%%%%%%%%%%%%%%%%%%%%%%%%%%%%%%%%%%%%

\begin{frame}
 \frametitle{Role of software development standards}
 \begin{itemize}
 \item Try to convey best practices in software development.
 \item Offer required framework for QA.
 \item Many examples: US DoD, ANSI, IEEE, BSI, NATO, etc.
 \item Engineers usually reluctant to adopt standards (bureaucracy).
 \end{itemize}
\end{frame}

%%%%%%%%%%%%%%%%%%%%%%%%%%%%%%%%%%%%%%%%%%%%%%%%%%%%%%%%%%%%%%

\begin{frame}
 \frametitle{ISO 9000}
 \begin{itemize}
 \item International body of standards for quality in any industry.
 \item ISO 9001 is the most extended version.
    \begin{itemize}
     \item Quality in process, development and maintenance of any product.
     \item Mainly concerned about high level procedures.
     \item Forces definition of specific metrics to measure quality.
     \item Involves periodic reviews to maintain ISO 9001 label award.
    \end{itemize}
 \item Involve definition of documentation templates, and documenting all processes.
 \item Does not ensure that standarized process reflects best practices in the field.
 \end{itemize}
\end{frame}

%%%%%%%%%%%%%%%%%%%%%%%%%%%%%%%%%%%%%%%%%%%%%%%%%%%%%%%%%%%%%%

\begin{frame}
 \frametitle{Documentation standards}
 \begin{itemize}
 \item 3 different categories:
 \item \textbf{Documentation process standards}: Establish good process for document production.
 \item \textbf{Documentation standards}: Define structure and presentation of documents.
 \item \textbf{Document interoperability standards}: Compatibility between document formats.
 \end{itemize}
\end{frame}

%%%%%%%%%%%%%%%%%%%%%%%%%%%%%%%%%%%%%%%%%%%%%%%%%%%%%%%%%%%%%%

\begin{frame}
 \frametitle{Documentation standards}
 \begin{itemize}
 \item Some examples
 \item \textbf{Document identification}: Unique names, easy tracking and identification.
 \item \textbf{Document structure}: Homogeneous structure for docs of similar character.
 \item \textbf{Document presentation}: Customized layout and ``look and feel'' (design and presentation).
 \item \textbf{Document update}: Tracking changes/updates/reviews.
 \end{itemize}
\end{frame}

%%%%%%%%%%%%%%%%%%%%%%%%%%%%%%%%%%%%%%%%%%%%%%%%%%%%%%%%%%%%%%

\section{Quality Planning}

%%%%%%%%%%%%%%%%%%%%%%%%%%%%%%%%%%%%%%%%%%%%%%%%%%%%%%%%%%%%%%

\begin{frame}
 \frametitle{Documentation standards}
 \begin{itemize}
 \item 5 different aspects:
 \item \textbf{Product overview}: Description, target market, user expectations.
 \item \textbf{Product plan}: Releases, updates, support, distribution, etc.
 \item \textbf{Process description}: Product development and management.
 \item \textbf{Quality goals}: Including identification of quality attributes to be tracked.
 \item \textbf{Risks and risk management}: Risks affecting quality and corrective actions.
 \end{itemize}
\end{frame}

%%%%%%%%%%%%%%%%%%%%%%%%%%%%%%%%%%%%%%%%%%%%%%%%%%%%%%%%%%%%%%

\begin{frame}
\frametitle{Software quality attributes}
\begin{center}
\begin{figure}
 \includegraphics[height=7.5cm]{figs/quality-attributes.pdf}
\end{figure}
\end{center}
\end{frame}

%%%%%%%%%%%%%%%%%%%%%%%%%%%%%%%%%%%%%%%%%%%%%%%%%%%%%%%%%%%%%%

\section{Quality Control}

%%%%%%%%%%%%%%%%%%%%%%%%%%%%%%%%%%%%%%%%%%%%%%%%%%%%%%%%%%%%%%

\begin{frame}
 \frametitle{Documentation standards}
 \begin{itemize}
 \item 3 components:
 \item \textbf{Quality reviews}: Periodic checks of conformity with quality standards define in the project.
 \item \textbf{Automated assessment}: Collecting quality information automatically (product and process).
 \item \textbf{Quality metrics}: Measure different quality aspects as an input for assessment.
 \end{itemize}
\end{frame}

%%%%%%%%%%%%%%%%%%%%%%%%%%%%%%%%%%%%%%%%%%%%%%%%%%%%%%%%%%%%%%

\begin{frame}
 \frametitle{Quality reviews}
 \begin{itemize}
 \item 3 components:
 \item We need a dedicated team of reviewers.
 \item Systematic methodology and approaches to examine key quality properties defined for the project.
 \begin{itemize}
  \item Design or program inspections.
  \item Progress reviews (management).
  \item Quality reviews: Technical analysis of components and assessment from metrics.
 \end{itemize}

 \end{itemize}
\end{frame}

%%%%%%%%%%%%%%%%%%%%%%%%%%%%%%%%%%%%%%%%%%%%%%%%%%%%%%%%%%%%%%

\begin{frame}
 \frametitle{Software measurement and metrics}
 \begin{itemize}
 \item \textit{Sofware metric}: Measurement linked to a software system, process or related documentation.
 \item Examples: SLOC, complexity, coupling, connascence, etc.
 \item Usually, the main problem is how to define meaningful metrics.
  \begin{itemize}
   \item We can measure many different aspects and elements.
   \item But do they describe the quality aspect we want to track?
   \item 2 types: \textit{Control} metrics (process) and \textit{predictor} metrics (product).
  \end{itemize}
  \item Key application: predicting possible problems for proactive management and control.
 \end{itemize}
\end{frame}

%%%%%%%%%%%%%%%%%%%%%%%%%%%%%%%%%%%%%%%%%%%%%%%%%%%%%%%%%%%%%%

\begin{frame}
\frametitle{Measurement procedure}
\begin{center}
\begin{figure}
 \includegraphics[height=7.5cm]{figs/quality-measurement.pdf}
\end{figure}
\end{center}
\end{frame}

%%%%%%%%%%%%%%%%%%%%%%%%%%%%%%%%%%%%%%%%%%%%%%%%%%%%%%%%%%%%%%

\section{Product quality}

%%%%%%%%%%%%%%%%%%%%%%%%%%%%%%%%%%%%%%%%%%%%%%%%%%%%%%%%%%%%%%

\begin{frame}
\frametitle{Product metrics}
\begin{center}
\begin{figure}
 \includegraphics[height=7cm]{figs/quality-product-metrics.pdf}
\end{figure}
\end{center}
\end{frame}

%%%%%%%%%%%%%%%%%%%%%%%%%%%%%%%%%%%%%%%%%%%%%%%%%%%%%%%%%%%%%%

\begin{frame}
 \frametitle{Product metrics}
 \begin{itemize}
 \item Static metrics:
 \begin{itemize}
  \item \textit{Fan-in/Fan-out}: Number of functions or methods that call some other
function or method (complexity, connascence).
  \item \textit{Length of code}: Size of program. SLOC (LOC or LLOC).
  \item \textit{Cyclomatic complexity}: Metric for control complexity of software.
  \item \textit{Length of identifiers}: Average length of identifiers used in a program.
Supposedly, the longer the ID, the better for readability and maintenance of code.
  \item \textit{Depth of conditional nesting}: Deeply nested statements are harder to be grasped.
  \item \textit{Fog index}: Average length of words and sentences in documents.
 \end{itemize}

 \end{itemize}
\end{frame}

%%%%%%%%%%%%%%%%%%%%%%%%%%%%%%%%%%%%%%%%%%%%%%%%%%%%%%%%%%%%%%

\begin{frame}
 \frametitle{Product metrics}
 \begin{itemize}
 \item Dynamic metrics:
 \begin{itemize}
  \item \textit{Depth of inheritance tree}: Number of discrete levels in the tree of
classes (OOP). The deeper the tree, the more complex the desgin.
  \item \textit{Method fan-in/fan-out}: Distinction between origin of calls to other
methods (from object or from external methods).
  \item \textit{Weighted methods per class}: Number of methods included in a class, weighted complexity
of every method. Too complex classes are difficult as for understanding and maintenance.
  \item \textit{Number of overriding operations}: Operations overridden in sub-classes. The higher this
number, the less appropriate may be the super-class.
 \end{itemize}

 \end{itemize}
\end{frame}

%%%%%%%%%%%%%%%%%%%%%%%%%%%%%%%%%%%%%%%%%%%%%%%%%%%%%%%%%%%%%%

\begin{frame}
 \frametitle{Metrics assessment}
 \begin{itemize}
 \item It is not enough to collect different metrics. We need to infer assessments and
conclusions from these results.
 \item Most of times, human interpretation of results is needed. 
 \item Value ranges drawn from multiple case examples may serve as guidelines.
 \item Usually, supported through integrated assessment frameworks (Business Intelligence).
  \begin{itemize}
   \item Organization of measurement areas and lists of metrics.
   \item Automated production of reports summarizing results.
   \item Support for elaboration of assessment reports and executive summaries.
   \item Other buit-in functionalities: data mining, data warehouse, statistical tools, etc.
  \end{itemize}

 \end{itemize}

\end{frame}

%%%%%%%%%%%%%%%%%%%%%%%%%%%%%%%%%%%%%%%%%%%%%%%%%%%%%%%%%%%%%%

\section{Process quality}

%%%%%%%%%%%%%%%%%%%%%%%%%%%%%%%%%%%%%%%%%%%%%%%%%%%%%%%%%%%%%%

\begin{frame}
 \frametitle{Cycle of process improvement}
\begin{center}
\begin{figure}
 \includegraphics[height=7cm]{figs/process-improvement-cycle.pdf}
\end{figure}
\end{center}
\end{frame}

%%%%%%%%%%%%%%%%%%%%%%%%%%%%%%%%%%%%%%%%%%%%%%%%%%%%%%%%%%%%%%

\begin{frame}
 \frametitle{Process improvement}
 \begin{itemize}
  \item Quality frameworks focused on improving the process, rather 
than just describing it.
  \item 3 stages:
  \begin{enumerate}
   \item \textit{Process measurement}: Compilation of metrics and indicators
related to software process (development and management).
   \item \textit{Process analysis}: Producing meaningful assessment from results,
identifying or anticipating problems that must be corrected.
   \item \textit{Process change}: Alter problematic elements or practices to align
the process to the quality goals of the project.
  \end{enumerate}
  \end{itemize}

\end{frame}

%%%%%%%%%%%%%%%%%%%%%%%%%%%%%%%%%%%%%%%%%%%%%%%%%%%%%%%%%%%%%%

\begin{frame}
 \frametitle{Link between process and product}
 \begin{itemize}
  \item Humphrey, in his book about process management (1988), states that
the same techniques applied to any kind of industry are valid for sofware
processes.
  \item However, a number of caveats are present in sofware projects.
  \begin{itemize}
   \item The design phase plays a key role in the whole process and results.
   \item Skills and experience of team leaders and experts affect quality of results.
   \item Different projects may need different methodologies and approaches.
   \item Great influence of 3 variables: cost, time frame and schedule.
  \end{itemize}

  \end{itemize}

\end{frame}

%%%%%%%%%%%%%%%%%%%%%%%%%%%%%%%%%%%%%%%%%%%%%%%%%%%%%%%%%%%%%%

\begin{frame}
 \frametitle{Process metrics}
 \begin{itemize}
  \item Process metrics should be aimed to facilitate process improvement.
  \item 3 types:
  \begin{itemize}
   \item Time to complete the process.
   \item Resources required to carry on the process.
   \item Number of occurrences of a particular event to be monitored.
  \end{itemize}

  \item The key question is to identify what we need to measure.
  \end{itemize}

\end{frame}

%%%%%%%%%%%%%%%%%%%%%%%%%%%%%%%%%%%%%%%%%%%%%%%%%%%%%%%%%%%%%%

\begin{frame}
 \frametitle{The GQM methodology}
 \begin{itemize}
  \item Frequently applied to develop customized process quality frameworks.
  \item Goal-oriented, top-down methodology to measure quality of software
processes and products.
  \item First developed at University of Maryland (1984) in collaboration with NASA.
  \item Hierarchical structure composed by 3 kinds of elements: goals, questions and
metrics.
  \end{itemize}

\end{frame}

%%%%%%%%%%%%%%%%%%%%%%%%%%%%%%%%%%%%%%%%%%%%%%%%%%%%%%%%%%%%%%

\begin{frame}
 \frametitle{GQM Hierarchy}
\begin{center}
\begin{figure}
 \includegraphics[height=6.5cm]{figs/gqm-abstract-example.pdf}
\end{figure}
\end{center}
\end{frame}

%%%%%%%%%%%%%%%%%%%%%%%%%%%%%%%%%%%%%%%%%%%%%%%%%%%%%%%%%%%%%%

\begin{frame}
 \frametitle{GQM implementation}
 \begin{itemize}
  \item We need to think at 3 different levels:
   \begin{itemize}
    \item \textbf{Conceptual} (goals): Goals must focus on an \textit{object},
is aimed to certain \textit{purposes}, pursue one or several \textit{quality focuses},
examine the object from one or different \textit{points of view} and it is
defined within a certain \textit{context}.
    \item \textbf{Operational} (questions): Break down one or several roles into
concrete issues to be inspected.
    \item \textbf{Quantitative} (metrics): Define measurable quantities or
logical aspects that providing information to offer precise answers and assessment in
response to the set of questions.
   \end{itemize}

  \end{itemize}

%%%%%%%%%%%%%%%%%%%%%%%%%%%%%%%%%%%%%%%%%%%%%%%%%%%%%%%%%%%%%%

\end{frame}

\begin{frame}
 \frametitle{GQM Example}
  \begin{center}
  \begin{figure}
    \includegraphics[height=5.75cm]{figs/gqm-example.pdf}
  \end{figure}
  \end{center}
\end{frame}

%%%%%%%%%%%%%%%%%%%%%%%%%%%%%%%%%%%%%%%%%%%%%%%%%%%%%%%%%%%%%%

\begin{frame}
 \frametitle{References}
 \begin{itemize}
  \item Humphrey, W. S. \textit{Characterizing the software process}.
IEEE Software, 5(2), 73-79. 1988.
  \item V. R. Basili, \textit{Quantitative Evaluation of Software Engineering Methodology}, 
Proceedings of the First Pan Pacific Computer Conference, Melbourne, 
Australia, Vol. 1 pp. 379-398, 1985.

  \end{itemize}

\end{frame}

%%%%%%%%%%%%%%%%%%%%%%%%%%%%%%%%%%%%%%%%%%%%%%%%%%%%%%%%%%%%%%

\end{document}